Fix a domain \(s\), a checked derivation \(\delta:A\Downarrow F\), and \(\mathbf M\in\mathcal C(\mathcal P)\).  We first prove
\[
 F(\mathcal P)=P_{M^s}(A)
\tag{10}
\]
by structural induction on the finite derivation tree in \(\mathrm{def:finite\text{-}extraction\text{-}calculus}\).

At an ordinary supplied-cell leaf, let \(T_a^s\) denote the finite tuple returned by regime \(a\).  By \(\mathrm{ass:typed\text{-}menu\text{-}realizability}\) and \(\mathrm{def:compatibility\text{-}fiber}\), for every complete tuple value \(u\),
\[
 \operatorname{Cell}(s,a,u)(\mathcal P)
 =P_a^s(T_a^s=u)
 =P_{M^s}(T_a^s=u).
\tag{11}
\]
The deterministic-postprocessing descriptor required by the typed grammar is also sound and adds no causal information.  If \(f\) is its declared map from the finite tuple alphabet to a finite output alphabet and \(v\) is an output value, its law is definitionally the pushforward of the already supplied tuple law.  Therefore
\[
 \begin{aligned}
 P_{a,f}^s(v)
 &=\sum_{u:f(u)=v}P_a^s(T_a^s=u)\\
 &=\sum_{u:f(u)=v}P_{M^s}(T_a^s=u) &&\text{by (11)}\\
 &=P_{M^s}\bigl(f(T_a^s)=v\bigr),
 \end{aligned}
\tag{12}
\]
where the last equality is finite additivity over the disjoint fibers of \(f\).  Taking \(f\) to be the identity recovers (11).  Moreover, evaluating \(f\) after the tuple is returned visits no new structural parent configuration, so the finite response carriers from \(\mathrm{def:district\text{-}response\text{-}profile}\) remain sufficient.

The contradiction and sure-event leaves have probabilities zero and one.  Suppose the induction claim holds for an event \(A_0\).  Finite complementarity gives
\[
 P_{M^s}(A_0^{\mathsf c})
 =1-P_{M^s}(A_0)
 =1-F_0(\mathcal P).
\tag{13}
\]
If it holds for a finite pairwise-disjoint family \((A_j:j\in J)\), finite additivity gives
\[
 P_{M^s}\!\left(\bigsqcup_{j\in J}A_j\right)
 =\sum_{j\in J}P_{M^s}(A_j)
 =\sum_{j\in J}F_j(\mathcal P).
\tag{14}
\]
Boolean normalization in the calculus produces precisely such a finite disjoint normal form, so (13)--(14) cover its Boolean constructors.  A marginalization constructor enumerates the finite partition \((A_x:x\in\mathcal X_W)\); hence (14) proves that case as well.

For a fixing constructor, the derivation contains an accepted fixing trace and type-correct substitutions for its leaves.  The trace preserves its represented probability by \(\mathrm{lem:finite\text{-}fixing\text{-}calculus\text{-}soundness}\); equations (11)--(14) and the induction hypotheses identify every substituted leaf.  Thus the input and output of the fixing constructor have the same probability.  An internal kernel quotient is accepted only after derivations of its numerator event \(A_0\cap H\) and context event \(H\), together with an independently accepted certificate that \(P_{M^s}(H)>0\).  Its induction step is consequently
\[
 \frac{F_{A_0\cap H}(\mathcal P)}{F_H(\mathcal P)}
 =\frac{P_{M^s}(A_0\cap H)}{P_{M^s}(H)}
 =P_{M^s}(A_0\mid H).
\tag{15}
\]
There is no unguarded division constructor.  These cases exhaust the checked derivation syntax, proving (10).

We next establish soundness of every accepted strict-positivity certificate without invoking a separation theorem.  Let \(x\) collect the finite response-law coordinates in the certificate, and write its declared constraint set as
\[
 h_i(x)=0\quad(1\le i\le r),
 \qquad g_j(x)\ge0\quad(1\le j\le t).
\tag{16}
\]
By the typing rule in \(\mathrm{def:finite\text{-}extraction\text{-}calculus}\), substitution of the coordinates induced by every model in the declared compatibility fiber satisfies (16).  An exact rational certificate is checked by reducing its claimed lower bound to a rational number \(\varepsilon>0\), which immediately proves the claim.  A linear face certificate supplies rational multipliers and an exactly checked identity of the form
\[
 p(x)-\varepsilon
 =\sum_{i=1}^{r}\alpha_i h_i(x)
  +\sum_{j=1}^{t}\lambda_j g_j(x),
 \qquad \varepsilon>0,\quad \lambda_j\ge0.
\tag{17}
\]
On (16), (17) implies \(p(x)\ge\varepsilon>0\).  The real-algebraic certificate supplies an exactly checked preordering identity
\[
 p(x)-\varepsilon
 =\sum_{i=1}^{r}\tau_i(x)h_i(x)
  +\sum_{J\subseteq\{1,\ldots,t\}}
    \sigma_J(x)\prod_{j\in J}g_j(x),
 \qquad \varepsilon>0,
\tag{18}
\]
where each displayed \(\sigma_J\) is itself checked as a finite sum of squares of rational polynomials.  Every summand in the second sum of (18) is nonnegative on (16), while the first sum vanishes; hence again \(p(x)\ge\varepsilon>0\).  Exact integer or rational coefficient comparison establishes (17) or (18) as a polynomial identity, so substitution cannot invalidate it.  These are all certificate constructors admitted by the definition.  Consequently each accepted positivity certificate is valid on its declared compatibility fiber, and in particular every division in (15) has a nonzero denominator.

Finally fix a district \(D\).  Define the finite nonfocal carrier
\[
 \Omega_{-D}=\prod_{D'\in\mathcal D(G),\,D'\ne D}\Omega_{D'},
\]
with the empty product interpreted as a singleton.  Finiteness follows from \(\mathrm{ass:recursive\text{-}compatibility}\) and \(\mathrm{def:district\text{-}response\text{-}profile}\).  Recursive evaluation of \(A\) therefore gives a total indicator
\[
 I_A:\Omega_D\times\Omega_{-D}\longrightarrow\{0,1\}.
\tag{19}
\]
By definition, \(\operatorname{SupportLocal}_D(A)\) checks the finite family of equalities
\[
 I_A(\omega_D,\omega_{-D})
 =I_A(\omega_D,\omega'_{-D})
\tag{20}
\]
for every \(\omega_D\in\Omega_D\) and every \(\omega_{-D},\omega'_{-D}\in\Omega_{-D}\).  Thus the procedure terminates.  If it accepts, (20) holds for every enumerated triple, which is exactly independence of the recursive indicator from the nonfocal profile.  Conversely, that independence makes every equality in (20) true, so the exhaustive checker accepts.

On acceptance, let \(\omega_{-D}^0\) be the checker's fixed reference profile.  Its constructed row is
\[
 \operatorname{row}_D(A)(\omega_D)
 =I_A(\omega_D,\omega_{-D}^0).
\tag{21}
\]
Combining (20) with (21) yields, for every \(\omega_D\in\Omega_D\) and \(\omega_{-D}\in\Omega_{-D}\),
\[
 \operatorname{row}_D(A)(\omega_D)
 =I_A(\omega_D,\omega_{-D}).
\]
If the checker rejects, the first failed equality in its finite order supplies an explicit triple witnessing dependence on the nonfocal profile.  This proves termination, the asserted if-and-only-if characterization, and correctness of the constructed row.
